% ECONOMICS_PROBLEM: {"id": "D83-640", "title": "AI-generated marketplace summaries", "jel_code": "D83", "source_id": "OP-0403"}
% !TeX program = xelatex
\documentclass[11pt,a4paper]{article}
\usepackage[margin=23mm]{geometry}
\usepackage{fontspec}
\setmainfont{Latin Modern Roman}
\usepackage{amsmath,amssymb,mathtools}
\usepackage{xcolor,enumitem,booktabs,longtable,array,xurl,hyperref}
\definecolor{muted}{HTML}{53636C}
\hypersetup{colorlinks=true,urlcolor=blue,linkcolor=blue}
\setlength{\parindent}{0pt}
\setlength{\parskip}{5pt}
\newcommand{\R}{\mathbb R}
\newcommand{\E}{\mathbb E}
\newcommand{\N}{\mathbb N}
\newcommand{\ind}{\mathbf 1}
\newcommand{\Var}{\operatorname{Var}}
\newcommand{\Cov}{\operatorname{Cov}}
\newcommand{\supp}{\operatorname{supp}}
\DeclareMathOperator*{\argmin}{arg\,min}
\DeclareMathOperator*{\argmax}{arg\,max}
\newcommand{\entrymeta}[3]{{\sffamily\small\color{muted}\textbf{\texttt{#1}}\quad|\quad #2\par #3\par}}
\newcommand{\Problemref}[1]{\texttt{#1}}
\newcommand{\JELref}[2]{\texttt{#2} (JEL \texttt{#1})}
\begin{document}
\section*{AI-generated marketplace summaries}
\textbf{Problem ID:} \texttt{D83-640}\quad \textbf{JEL:} \texttt{D83}\par
% BEGIN ECONOMICS STATEMENT
\label{problem:OP-0403}
{\sffamily\small\color{muted}Original subject group: Digital, platform and data economics\par}
\label{definitions:problem:30007685}
\textbf{Audit classification:} Published research agenda. The source motivates average and distributional effects of AI summaries. Product-cluster assignment does not alone justify scalar potential outcomes for multi-product viewers.
\subsection*{Definitions and precise research target}
Consider product reviews on an e-commerce platform and the population of users viewing eligible products during a prespecified six-month experiment. Randomly assign each product-market cluster to \(a\in\{0,1\}\), where \(a=1\) displays an AI-generated review summary and \(a=0\) displays the existing review interface. Hold access to full reviews available in both regimes. Let \(Y_i(a)\) be buyer \(i\)'s independently measured money-valued postpurchase satisfaction minus price and search cost, assigning zero gross satisfaction to the outside option. Include returns and unmet expectations in the satisfaction measurement. Define \(\tau=E[Y_i(1)-Y_i(0)]\) and \(\tau_g=E[Y_i(1)-Y_i(0)\mid G_i=g]\), where \(G_i\) is a prespecified buyer group and expectation averages over the experiment's eligible-user population. Compare summaries with audited review evidence to measure omitted adverse information and factual errors. Account for seller adaptation and review manipulation by separately estimating immediate effects and effects after the six-month deployment. Specify exposure mappings for buyers who view products in multiple clusters; otherwise report interference sensitivity bounds. Determine whether summaries improve match value and whether gains concentrate among particular users or sellers. Engagement is a secondary outcome, because more clicks alone do not define better matches.

\textbf{Definitions, assumptions and mathematical qualifications.} The eligible population includes all viewers, not just purchasers. A summary policy includes model version, prompt, review cutoff, refresh rule and full-review access. With several clusters viewed by one person, \(Y_i(z)\) is indexed by the full assignment vector \(z\); a scalar \(Y_i(a)\) is valid only under no relevant cross-cluster interference or a prespecified exposure mapping. Define satisfaction using one monetary elicitation or utility model across arms; outside-option gross value and purchase expenditure are zero but incurred search cost remains. Separate immediate outcomes from six-month policy outcomes incorporating seller/review responses.
\subsection*{Variants and assumption changes}
\begin{enumerate}
\item \textbf{Editorial, fixed review corpus.} Freeze eligible reviews and seller offers and randomize summaries, identifying a direct information effect under no cross-cluster spillovers.
\item \textbf{Editorial, adaptive review ecosystem.} Randomize geographically complete markets, permit new reviews and seller responses, and include factual-error, omitted-adverse-information and group-welfare outcomes after six months.
\end{enumerate}
\subsection*{References and source audit}
\textbf{Checked source.} 2025 conference draft, Section 1.1, fourth research direction, PDF p. 5. Full primary text retrieved. \url{https://conference.nber.org/confer/2025/DTs25/farronato.pdf}.
\begin{enumerate}\raggedright
\item\hypertarget{definitions:source:30007685:1}{} Chiara Farronato. 2025. \emph{Digital Platforms as Information Aggregators: Implications for their Design and Regulation}. 2025 conference draft, Section 1.1, fourth research direction, PDF p. 5.
\url{https://conference.nber.org/confer/2025/DTs25/farronato.pdf}.
\end{enumerate}

\subsection*{Common conventions: Source mathematical conventions}

These conventions apply to each entry unless explicitly replaced.

\textbf{Models and identification.} A primitive model \(m\) specifies all probability laws, feasible choices, information, equilibrium and measurement rules used by its target. The admissible class \(\mathcal M\) is fixed independently of outcomes. If \(P_O(m)\) is its observable probability law and \(\tau(m)\) the target, its identified set at observable law \(P\) is
\[
 \mathcal I_\tau(P)=\{\tau(m):m\in\mathcal M,\ P_O(m)=P\}.
\]
Point identification means that this set is a singleton. An empty set signals incompatibility of data and assumptions. Coordinate bounds need not describe the joint set. Bounds are sharp only if every retained target is attainable or arbitrarily approachable by compatible models. Finite bounds require explicit restrictions on unobserved quantities; unrestricted confounding, hidden wealth, extrapolation or arbitrary equilibrium selection can yield uninformative sets.

\textbf{Causal interventions.} A potential outcome \(Y_i(p)\) is the outcome under a complete policy regime \(p\), including timing, financing, information and market spillovers. The target population and units are fixed before treatment. With interference, use the complete assignment vector or a justified exposure mapping; individual no-interference notation alone cannot identify an economy-wide intervention. A difference of mean potential outcomes does not require the cross-world joint distribution. Conditioning on post-treatment migration, survival, enrollment or employment changes the estimand and needs separate justification.

\textbf{Statistical statements.} An estimator, confidence set, forecast or test specifies a sampling law, model class and loss/coverage criterion. Uniform \((1-\alpha)\)-coverage means \(\inf_{m\in\mathcal M}\Pr_m\{\tau(m)\in C_n\}\geq1-\alpha\), or an explicitly asymptotic version. The whole parameter space trivially has coverage, so informativeness requires a stated width, risk or power objective. A forecasting improvement is relative to a named benchmark, information set and evaluation population.

\textbf{Equilibrium and optimization.} An equilibrium comprises feasible plans, optimality, beliefs where needed, budget constraints and all market/resource clearing. The primitives or a stated benchmark define each correspondence \(\mathcal E(p;m)\). Nonemptiness is assumed or proved, never inferred just from compact policy sets. A selected-equilibrium target uses a declared selection rule; a robust target quantifies over all permitted equilibria. Use suprema or infima unless attainment is established. An \(\varepsilon\)-optimal policy is feasible and within \(\varepsilon\) of the finite optimum value. Report an empty feasible set separately.

\textbf{Welfare and accounting.} Normative utility, interpersonal normalization, social weights, numeraire, discounting and the use of public receipts are primitives. Behavioral utility need not coincide with the welfare criterion. Household welfare includes ownership income and rents once; adding profits again to compensating variation generally double-counts them. A monetary equivalent is defined by an explicit transfer and price convention, with monotonicity and range conditions. Average effects, mechanism contributions, transfers and welfare losses are different targets. Infinite sums require convergence and dynamic feasible plans require terminal or no-Ponzi conditions.

\textbf{Source versions.} A working-paper date, revision date and journal publication year are distinguished. A DOI for a journal version is not automatically the DOI of an earlier manuscript. Locators refer to the stated version and printed pages where specified. A metadata-only check does not verify a claim about a full-text conclusion. Every entry's source subsection narrows the provenance claim accordingly.
% END ECONOMICS STATEMENT
\end{document}
